\scp{Languages in linguistic Wallacea}{15-01}
If we apply a similar concept to language,
we start out far to the north
and west of the Blust line with languages that are clearly
and indisputably Austronesian in their morphosyntax,
phonology and lexicon.\footnote{
West and north of the Blust line
there are known areas of contact,
such as the influence of Mon-Khmer languages in Aceh
and in Chamic languages.
Unusual features in the \ili{Batak} languages of Sumatra are sometimes
assumed to reflect some sort of substrate influence,
as are the Agta languages of the Philippines,
and the Sama/Orang Asli languages of Malaysia.}
Many Austronesian
languages in those regions are typologically VSO or SVO in their
unmarked word order in transitive clauses with NP core arguments
and elaborate voice systems.
By the time they reach what is now central Indonesia the dominant
typology is SVO with prepositions,
and with post-posed possessors,
as in \ili{Malay} \it{rumah=ku} or \it{rumah saya} ‘house \tsc{1sg}’ = ‘my house’.
Body parts tend to be possessed in the same way as other nouns,
as with \ili{Malay} \it{taŋan=ku} or \it{taŋan saya} ‘hand \tsc{1sg}’
= ‘my hand’, and \it{taŋan} can stand alone without a possessor.
These languages overwhelmingly tend to have their standard negation
before the verb (pre-predicate).

Far to the \ili{east},
in the region of the island of New Guinea,
there are hundreds of languages that are clearly not Austronesian,
but are affiliated with dozens of unrelated language families
popularly lumped together under the label of ‘Papuan’, basically meaning ‘non-Austronesian’.
The dominant typology in the \tsc{North Halmahera}
and the \tsc{Timor-\ili{Alor}-Pantar} languages is verb-final SOV, with pre-posed possessors.
These languages tend to have their standard negation after the
verb, which commonly puts it at the end of the clause.
They often have postpositions.
In their possessive systems they tend to distinguish between
alienable and inalienable noun classes or possessive constructions.
And they often have a human/non-human distinction in their third person pronominal systems.

In the middle between the two extremes,
there are some Papuan languages still spoken in parts of our target region.
There is some direct evidence,
and an abundance of trace evidence,
that Papuan languages were once spoken much more widely throughout
the region than they are today (see \srf{02-01} and \srf{02-03-01}).
Those that are still spoken have mostly maintained their SOV
typology with pre-posed possessors.
A few have shifted to SVO typology associated with Austronesian languages.
And some have developed an inclusive-exclusive distinction in
\tsc{1pl} pronouns associated with Austronesian languages.
Austronesian loanwards in Papuan languages are widespread,
and in some Papuan languages Austronesian loans may dominate
whole segments of certain semantic fields,
such as numbers or colours.\footnote{
		Grimes found this was the
		case in a variety of \ili{Fataluku} (a Papuan TAP language spoken on
		the eastern tip of Timor) from which he collected data in November
		1999 (see also \cite[26]{Heston2015}, \cite{UsherSchapper2018,UsherSchapper2022}).} 
\ili{As} discussed in \srf{02-03-01},
most Austronesian languages in the target region can be characterised
as having preposed possessors more commonly associated with the
SOV typology of Papuan languages,
not commonly associated with the prepositional SVO typology of
the Austronesian languages and not generally found in Austronesian languages west
and north of the Blust line.
But preposed possessors are also not found in much of the \tsc{\ili{Bima}-Lembata} subgroup (see \frf{02-16}).
Large numbers of Austronesian languages have shifted to clause-final
negation -- another feature associated with the SOV typology of
Papuan languages, but not commonly associated with the SVO typology of most Austronesian languages (see \srf{02-03-01-02}).
Some pockets of Austronesian languages show a distinction between
alienable and inalienable nouns -- another feature commonly found in Papuan languages (see \srf{02-03-01-01}).
The same is true for the appearance in pockets of languages of
a distinction between human/non-human marking in third person
pronominal systems (see \srf{07-04-03}).

These typological features associated with Papuan languages are
found widely, but not distributed evenly or uniformly throughout the region.
\ili{As} areal features resulting from contact they are not useful
for linguistic subgrouping at higher levels.

Each of these structural
and typological changes individually imply intense contact
and bilingualism over a long period of time.
Collectively these features mean we are talking about a group
of a very different type of language than those found to the
west and north of the Blust line.
One would not be wrong in thinking of these languages as hybrids
with many different combinations
and degrees of language mixing.
But those structural features mentioned are typological due to
contact, not genealogical due to inheritance from a common Austronesian
parent language or linkage.

Add to these some of the changes noted in basic vocabulary (e.g.
the shift to ‘leaf of head’ to ‘hair’; see \srf{02-03-01-05}),
basic metaphors of life (e.g.
the shift from a secondary sense of ‘liver’ to a secondary sense
of ‘inside’ to talk about ‘seat of emotions,
character and social relations’; see \srf{02-03-01-06}) away
from seemingly stable concepts in Austronesian languages west
and north of the Blust line to accommodate common Papuan ways
of talking about these things.
Once again, the distribution is not even across the target region,
not all languages have them,
and some features are hardly found in some subgroups,
or not at all.
For example: 

\begin{itemize}
\item	East of the Blust line,
in western Flores and \tsc{Sumba-Havu} languages, preposed possessors
associated with Papuan languages do not appear to be found (see \frf{02-16}).
In fact a number of features commonly associated with this Wallacean
region are not found in these languages.
Yet Papuan-sourced {\#}muku ‘banana’ made it as far west as \ili{Manggarai}
in western Flores (\frf{02-20}).
\ili{Dhao} and \ili{Havu} use ‘leaf of head’ to talk about ‘hair’,
and ‘inside’ to talk about ‘seat of emotions and character’.
So some traces of contact with Papuan languages are evident in
languages that do not have preposed possessors.
The transition in languages from Java to Bali to Lombok to \ili{Sumbawa}
to western Flores is relatively gradual,
with a boundary in the middle of \ili{Sumbawa} separating the western languages from \tsc{\ili{Bima}-Lembata}.

\item	Near the northern parts of our target region,
just west of the Blust line,
in \ili{Banggai} (Sulawesi) and \ili{Taliabu}, preposed possessors are also found.
But those languages link westward with the \tsc{Celebic} languages
of Sulawesi in their historical phonologies,
not eastward with other Wallacean subgroups that are associated with preposed possessors.
(See \crf{ch:Tal}.) 
\item	Travelling across Flores,
there is a pocket of languages in \ili{Central} Flores that appear
to have very little morphology (Elias 2018, 2020; McWhorter 2019).
These form a node within our \tsc{\ili{Bima}-Lembata} subgroup.
Does that reflect substrate influence? Some sort of incursion?
Language shift? Imperfect second language acquisition by adults?
Or simply ``normal'' processes of language change?

\item	Travelling across Timor,
there is a subgroup in \tsc{\ili{Central} Timor} that has a different
phonological history than the surrounding \tsc{Timor-Babar} languages.
\tsc{Timor-Babar} itself has huge variation within it.
(See chapters \ref{ch:TB} and \ref{ch:CT}).
Both are in direct
and on-going contact with Papuan languages.
Is the diversity from substrate influence? Incursion? Innovation?
All of the above?

\item	Within neighbouring subgroups or nodes,
			contact-induced change may characterise one branch quite strongly,
			and another branch to a significantly lesser degree.
			For example, \tsc{\ili{Ambon}-Seram} languages have subject indexing
			on verbs, multiple inflectional paradigms with different morphophonemic
			changes, alienable-inalienable possession, and human/non-human distinction
			in third person pronouns (see \srf{07-04-03}).
			Nearby \tsc{Sula-Buru} languages have very little to none of these.
			Yet both have preposed possessors,
			both have some languages with clause final negation,
			both have languages with a secondary sense of ‘inside’ as ‘seat
			of emotion and character’, and neither subgroup has languages with reflexes
			of Papuan {\#}muku ‘banana’.
			There are very consistent patterns in the historical phonology
			and words in the lexicon that distinguish the two subgroups -- along
			with a few words that are shared across both subgroups.
			Do the differences in \tsc{\ili{Ambon}-Seram} languages reflect substrate
			influence? Or some sort of incursion? Which group is more conservative?
			Which did more innovating? Regardless of the answers to those
			questions, the two subgroups clearly have different phonological
			histories (chapters \ref{ch:CM}, \ref{ch:SB}, \ref{ch:AS}).
\item	When we reach \ili{Aru} or the Bomberai peninsula of west Papua,
			people who are familiar with Austronesian languages in the western
			parts of Indonesia are simply shocked at first encounter that
			these languages could be considered Austronesian.
			The Papuan influence is so strong on the surface,
			that when looking at whole sentences
			and texts, many of these languages do not initially look,
			or sound or act very “Austronesian”.
			To get a feel for this,
			readers who are familiar with \ili{Malay},
			\ili{Balinese}, \ili{Tagalog}, or Cebuano should skim the \ili{Dobel} examples provided in \citet{Hughes2000}.
			One could rightly ask,
			“Are these really Papuan languages in Austronesian disguise?
			Or maybe Austronesian languages in Papuan disguise?”\footnote{
					Questions
					posed by Mark Donohue in 2011,
					and reiterated in \citet{DonohueDenham2020}.}
			On what basis
			do we consider them to be Austronesian? Digging below the surface,
			we find there are consistent
			and systematic patterns in the phonology,
			the lexicon, and the morphosyntax that are very Austronesian,
			and are too pervasive,
			systematic and deep-seated to be accounted for by contact.
			Hence we say they are Austronesian languages.
			But they are also not “Austronesian” in the same way as languages
			in western Indonesia or the Philippines.
\item	There is a corridor of Austronesian languages stretching from
			\ili{Aru} to South Halmahera,
			and including the Bird’s Neck of Papua.
			Within this corridor a number of languages have monosyllables
			for many words which are reconstructed in PMP with disyllables or trisyllables.
			This areal tendency is found in some languages in three different
			subgroups (\tsc{Aru}, \tsc{Seram-Tanimbar-Bomberai}, SHWNG).
			That means the already limited data is even more restricted for
			discovering what are the systematic sound correspondence in these
			languages for vowels in final syllables (which have subgrouping
			value in \tsc{Seram-Tanimbar-Bomberai} languages),
			vowel-glide sequences such as *-ay,
			*-aw, *-uy, and final syllables which are often lost altogether.
\item	In \crf{ch:STB},
			it was noted that the \tsc{Bomberai Tip}
			and \tsc{East Bomberai} languages continue to be in regular intense
			contact with Papuan languages.
			It was also noted there that these languages are so heavily Papuanised
			that finding enough Austronesian cognates to establish regular
			sound correspondences is often difficult.
			And this is further complicated in these areas of traditional
			sago-prominant diets, by having words such as ‘rice’ taken from other Austronesian
			languages such as \ili{Geser},
			which has functioned for centuries as a lingua franca in that region.
			Once again it is fair to ask,
			on what basis do we consider them to be Austronesian?
			And again, digging below the surface,
			we find there are consistent
			and systematic patterns in the phonology,
			the lexicon, and the morphosyntax that are very Austronesian,
			and are too pervasive
			and systematic in the languages to be accounted for by contact.	
\end{itemize}

Just as biogeographical Wallacea is a zone of transition between
the flora and fauna of Asia (Sundaland)
and the flora and fauna of Australia-New Guinea (Sahulland),
so also our target area in Linguistic Wallacea is a zone of transition
between Island South\ili{east} Asian (Austronesian) languages,
and Melanesian (Papuan) languages.
So at the western edge of the zone of transition we find the
languages in western Flores
and Sumba-\ili{Havu} show only a few features that are associated with Papuan languages.
But there are some.
At the other extreme,
at the eastern edge of the zone of transition some of the \tsc{Aru}
languages, along with the Bomberai Tip
and East Bomberai languages,
have such heavy influence from a Papuan substrate that they are
barely recognisable as Austronesian.